% ---------------------------------------------------------------------
\section{Request pipeline and perfect-information optima}\label{sec:pipeline}
% ---------------------------------------------------------------------

\subsection{Station-level request pipeline}\label{app:admission}

The arrival cap and admission control belong to the request setting of the station operator and not to the controller.
For every test replay they are computed once, before any controller runs, and the resulting admitted requests $\Ereq_i$ are shared by every compared method and by both optima.
The wrapper writes the admitted requirement into EV2Gym's EV object at arrival, so every controller that reads EV2Gym's EV objects pursues the same target.
For arrivals inside its horizon, the HiGHS eMPC uses the shared entries.
Conformal passes re-certify the shared requirement with their bound $\hat d_i^{\sm{lo}}$, so their requirement is never larger, and passes without admission control keep the arrival-capped request.
In every pass, shortfall is measured against the shared admitted requirement.

Arrivals are processed in time order, with ties broken by the charger index.
An arriving EV $j$ enters with its exogenous arrival energy, desired energy, departure step and EV parameters, and the cap of \Eq{M-eq:cap} gives its first requirement with a certified threshold table.
Admission control then solves the feasibility LP of \Eq{eq:admission} at the arrival step $t_0$ over the EVs admitted earlier and still connected, together with EV $j$:
\begin{equation}\label{eq:admission}
\begin{aligned}
&E_i(t+1)=E_i(t)+\Delta t\,p_i(t),\qquad E_i^{\min}\le E_i(t+1)\le E_i^{\sm{cap}},\qquad E_i(d_i+1)\ge E_i^{\sm{tgt}},\\
&-\Pdis_i\le p_i(t)\le\Pnom_i,\qquad p_i(t)\le\varrho_i^{\sm{in}}\,(E_i^{\sm{cap}}-E_i(t))\ (t>t_0),\qquad p_i(t_0)\in[\pmn_i,\bar P_i],\\
&-\hat C_{\sm{rev}}(t)\le\textstyle\sum_ip_i(t)\le\hat C(t)-n(t_0)\varepsilon_q .
\end{aligned}
\end{equation}
The energies of the earlier EVs at $t_0$ are read from a nominal plan and not from the trajectory of any controller.
After each admission decision, an elastic version of the same LP updates this plan so that the targets are reached as early as possible.
The limits $\hat C$ and $\hat C_{\sm{rev}}$ use the measured values at $t_0$, the noisy forecasts for the next $H_f$ steps and the last forecast afterwards.
Its slope $\varrho_i^{\sm{in}}$ is the largest value with $\bar P_i(E)\ge\min\{\Pnom_i,\varrho_i^{\sm{in}}(E_i^{\sm{cap}}-E)\}$ for all $E$, so the linear taper model is an inner approximation and every LP-feasible schedule is physically realisable.

Inside the LP, the new EV has the target $E_j^{\sm{tgt}}=\Ereq_j$, and an earlier EV has the smaller of $\Ereq_i$ and $\Psi_i^{\sm{LP}}(E_i(t_0))-10^{-6}$ as its target.
Here $\Psi_i^{\sm{LP}}$ is the largest departure energy reachable in the LP's taper model, and this target is used only inside the admission LP.
With full-battery requests as literal targets, the LP would be infeasible whenever such an EV is connected, because the linear taper model never reaches full charge in finitely many steps.
Feasibility is monotone in the new request, so one LP that maximises $E_j(d_j+1)$ gives the largest admissible request.
Requests admitted earlier are never reduced, and the reduced request of EV $j$ passes through the cap routine again, so A4 holds for every EV and every method.
If the LP is infeasible even for the lowest request, the request is set to the arrival energy.

Each decision uses only information available at $t_0$, so a station operator can run the pipeline online.
If its LP is feasible, an open-loop schedule exists that meets all targets from the nominal-plan energies under the forecast limits, provided that no further EV arrives.
The pipeline does not guarantee A5 for any controller, because realised energies may lie below the nominal plan, forecasts are noisy and later arrivals add load.
It therefore does not enter Theorem~\ref{M-thm:invariance}, and the evaluation passes without admission control measure its effect.

\subsection{Linear programs of the optima}\label{app:oracle}

Both perfect-information LPs know all sessions of the episode with their arrival and departure steps, arrival energies and admitted requirements, together with all prices and the realised limits $C(t)$ and $\Crev(t)$.
For every session $i$ and step $t\in\mathcal T_i=\{t_i^{\sm{arr}},\dots,d_i\}$ the variables are $p_i^{+}(t),p_i^{-}(t)\ge0$ and $E_i(t)$, together with a shortfall slack $z_i\ge0$ and, in the soft LP, overload slacks $o^{\pm}(t)\ge0$.
With the rounding allowance $G_i^{\max}=0.01\,(d_i-t_i^{\sm{arr}}+1)$ kWh, the constraints are
\begin{equation}\label{eq:lp}
\begin{aligned}
&E_i(t_i^{\sm{arr}})=E_i^{\sm{arr}},\qquad E_i(t+1)=E_i(t)+\Delta t\,\big(p_i^{+}(t)-p_i^{-}(t)\big),\qquad E_i^{\min}-G_i^{\max}\le E_i(t+1)\le E_i^{\sm{cap}},\\
&p_i^{+}(t)\le\Pnom_i,\qquad p_i^{-}(t)\le\Pdis_i,\qquad p_i^{+}(t_i^{\sm{arr}})\le\bar P_i(E_i^{\sm{arr}}),\qquad p_i^{+}(t)\le\varrho_i^{\sm{out}}\big(E_i^{\sm{cap}}-E_i(t)\big)\ (t>t_i^{\sm{arr}}),\\
&E_i(d_i+1)+z_i\ge\Ereq_i-G_i^{\max},\qquad -\Crev(t)-o^{-}(t)\le\textstyle\sum_{i}\big(p_i^{+}(t)-p_i^{-}(t)\big)\le C(t)+o^{+}(t),
\end{aligned}
\end{equation}
where the hard LP fixes $o^{\pm}=0$.
The profit of a schedule is $\mathrm{Pr}(p)=\sum_t\Delta t\sum_i\big(c^{\sm{dis}}p_i^{-}-c^{\sm{ch}}p_i^{+}-c^{\sm{thr}}(p_i^{+}+p_i^{-})\big)$, with the throughput cost $c^{\sm{thr}}$ equal to zero in all reported runs.

\begin{lemma}[Taper rows]\label{lem:taperrows}
Let every EV satisfy the conditions of Lemma~\ref{lem:a3}, and let $\varrho_{\kappa,i}=(1-e^{-\kappa_i})/\Delta t$.
Then $\bar P_i(E)\le\min\{\Pnom_i,\varrho_{\kappa,i}(E_i^{\sm{cap}}-E)\}$ for every $E\in[E^{\min},E^{\sm{cap}}]$.
Equality holds except on the one-step transition band $\varsigma\in(\varsigma_i^{\sm{tr}}-\delta_i,\varsigma_i^{\sm{tr}})$, and $\sup_E\bar P_i(E)/(E_i^{\sm{cap}}-E)=\varrho_{\kappa,i}$.
\end{lemma}

\begin{proof}
In state-of-charge units, $\bar P\Delta t/E^{\sm{cap}}=\varphi/E^{\sm{cap}}-\varsigma$ by \Eq{eq:taper}, and for $\varsigma\ge\varsigma^{\sm{tr}}$ this gain is $\min\{\delta,(1-\varsigma)(1-e^{-\kappa})\}$ and equals the right-hand side.
The constant-current region $\varsigma+\delta\le\varsigma^{\sm{tr}}$ has the gain $\delta$, and the right-hand side equals $\delta$ if and only if $\delta\le(1-\varsigma)(1-e^{-\kappa})$.
Since $1-\varsigma\ge1-\varsigma^{\sm{tr}}=m\delta/\kappa$, it suffices that $\kappa\le m(1-e^{-\kappa})$, and with $\kappa=mz$ this is $z\le1-e^{-mz}$, that is, $\kappa\le\kappa^{*}(m)$.
On the transition band, let $y=\varsigma^{\sm{tr}}-\varsigma\in(0,\delta)$, so that $1-\varphi/E^{\sm{cap}}\ge(1-\varsigma^{\sm{tr}})e^{-m(\delta-y)/(1-\varsigma^{\sm{tr}})}$.
With $e^{v}\ge1+v$ and $m\ge1$, the right-hand side of this bound is at least $e^{-\kappa}(1-\varsigma^{\sm{tr}}+my)\ge e^{-\kappa}(1-\varsigma)$.
The gain is therefore at most $(1-\varsigma)(1-e^{-\kappa})$, and it is at most $\delta$ anyway, so the inequality holds everywhere, and the supremum follows from equality in the upper region.
\end{proof}

The code uses $\varrho_i^{\sm{out}}=1.01\,\varrho_{\kappa,i}$, a valid outer row that is 1\% looser than exact.
A check script released with the code confirms the inequality to $1.7\times10^{-13}$ kW and the equality off the band on 320 configurations.
With these rows the LPs model EV2Gym's charging limit up to the step discretisation of the transition band and the 1\% margin.
The remaining relaxations are the rounding allowance and simultaneous charging and discharging, and the latter is cost- and energy-neutral, and hence exact, when $\eta=1$, $c^{\sm{dis}}=c^{\sm{ch}}$ and $c^{\sm{thr}}=0$, as in all reported runs.

\subsection{Upper bounds}\label{app:bounds}

The hard LP is solved lexicographically with HiGHS \cite{huangfu2018parallelizing}.
Stage 1 computes the least total shortfall $z^{*}=\min\sum_iz_i$ with $G^{\max}=0$, and stage 2 maximises $\mathrm{Pr}(p)$ with the allowance and $\sum_iz_i\le z^{*}+10^{-6}$, giving $J^{\sm{LP,ceil}}$.
The same stage 2 with $G^{\max}=0$ gives $J^{\sm{LP}}$, reported only in \Sec{sec:results-supp}.
In contrast, the soft LP maximises $\mathrm{Pr}(p)-w_E\sum_iz_i-w_P\Delta t\sum_t(o^{+}(t)+o^{-}(t))$ subject to \Eq{eq:lp} in a single stage, giving $J^{\sm{pen,LP}}$.

\begin{proposition}[Hard optimum]\label{prop:oracle}
Let the configuration of Lemma~\ref{lem:a3} hold, and let $\pi$ be any controller whose reported powers satisfy \Eq{M-eq:transformer} at every step and whose realised total shortfall against the admitted requirements is at most $z^{*}$.
Then $J^{\pi}\le J^{\sm{LP,ceil}}$.
\end{proposition}

\begin{proof}
Let $\tilde p_i(t)$ be the reported powers and $E_i^{\sm{nom}}(t)=E_i^{\sm{arr}}+\Delta t\sum_{s<t}\tilde p_i(s)$ the energy that the billed powers alone would produce.
EV2Gym bills exactly $\tilde p$, so $J^{\pi}$ equals $\mathrm{Pr}$ at $p^{\pm}=\tilde p^{\pm}$.
The rounding adds less than 0.01 kWh per active step, so $E_i^{\sm{real}}(t)-G_i^{\max}\le E_i^{\sm{nom}}(t)\le E_i^{\sm{real}}(t)\le E_i^{\sm{cap}}$ along the trajectory.
Both the lowered energy floor and the lowered departure row therefore hold for $E^{\sm{nom}}$ with $z_i=(\Ereq_i-E_i^{\sm{real}}(d_i+1))^{+}$, and the power rows hold because $\tilde p_i\le\bar P_i(E_i^{\sm{real}})\le\Pnom_i$ and $\tilde p_i\ge-\Pdis_i$.
Moreover, the first-step row holds because $E^{\sm{nom}}(t_i^{\sm{arr}})=E_i^{\sm{arr}}$.
For the taper rows, EV2Gym never delivers more than $\varphi(E)-E$, so Lemma~\ref{lem:taperrows} gives $\tilde p_i(t)\le\varrho_{\kappa,i}(E^{\sm{cap}}-E_i^{\sm{real}}(t))\le\varrho_i^{\sm{out}}(E^{\sm{cap}}-E_i^{\sm{nom}}(t))$.
The transformer rows are the assumption, and $\sum_iz_i\le z^{*}$, so the trajectory is feasible for stage 2 and its objective equals $J^{\pi}$.
\end{proof}

\begin{proposition}[Soft-penalty optimum]\label{prop:softoracle}
Let the configuration of Lemma~\ref{lem:a3} hold, and let commands be clipped to the physical limits of \Eq{M-eq:limits}.
For every simulated trajectory on a replay, with shortfall measured against the shared admitted requirement, $J^{\sm{pen}}(\pi)\le J^{\sm{pen,LP}}$.
\end{proposition}

\begin{proof}
Take $p^{\pm}=\tilde p^{\pm}$, $E=E^{\sm{nom}}$, $z_i=\Delta E_i^{\sm{sh}}$ and $o^{\pm}(t)$ equal to the realised overloads in either direction.
The energy, power, first-step and taper rows hold as in the proof of Proposition~\ref{prop:oracle}, because they use only the physics of EV2Gym and not any constraint that the controller respects.
For the departure row, $E^{\sm{nom}}(d_i+1)+z_i\ge E^{\sm{real}}(d_i+1)-G_i^{\max}+z_i\ge\Ereq_i-G_i^{\max}$, and the transformer rows hold by the definition of $o^{\pm}$.
At this point the objective equals $J^{\sm{pen}}(\pi)$, with the same billed powers, shortfall and overload energy.
\end{proof}

The bound is a property of the measurement, not of the controller, so it covers controllers with shortfall or overload, the conformal and no-admission passes and the rule-based and model-predictive baselines.
Every pass in the wrapper applies the physical clip.
In the reported runs, none of the 153,100 episodes exceeded the soft-penalty optimum by more than $10^{-6}$\,\texteuro{} (analysis load report).
A unit test released with the code checks $J^{\sm{pen}}\le J^{\sm{pen,LP}}$ for all rule-based methods on small replays, including random actions.
Inside EV2Gym, a controller can exceed $J^{\sm{LP}}$ through the unbilled rounding gain, so the main text uses $J^{\sm{LP,ceil}}$ for the cost of the guarantee and reports $J^{\sm{LP}}$ only here.
